\documentclass[12pt,a4paper]{article}
\usepackage[utf8]{inputenc}
\usepackage[T1]{fontenc}
\usepackage{amsmath,amssymb,graphicx,booktabs,hyperref,natbib}
\usepackage[margin=2.5cm]{geometry}

\begin{document}

\title{\bf A non-monotonic dark energy equation of state\\
from cosmic structure formation}

\author{Anonymous}
\date{June 2026}

\maketitle

\begin{abstract}
If cosmic structure formation occupies a finite reservoir of vacuum quantum modes, the dark energy equation of state $w(z)$ inherits the shape of the cosmic star formation history. We solve the resulting model self-consistently---iterating the coupled Friedmann and mode-occupation equations---and find that the converged solution predicts a non-monotonic $w(z)$: $w_0=-0.80$, rising to a peak of $w\approx-0.31$ at $z\approx1.5$, then returning toward $-1$ at higher redshift. The Hubble parameter deviation from $\Lambda$CDM is approximately $2\%$ at all measured redshifts. Against DESI DR2 baryon acoustic oscillation distances with the sound horizon as a free parameter, the model yields $\chi^2=19.1$ compared to $\chi^2=15.0$ for $\Lambda$CDM ($\Delta{\rm AIC}=+8.1$). The model is not excluded, though not statistically preferred. Its distinguishing feature---a peaked equation of state---cannot be produced by the CPL parameterization or by quintessence with simple potentials for any choice of parameters. DESI DR3 and the Euclid mission, with projected binned $w(z)$ precision of $\sigma_w\sim0.1$--$0.2$ per $\Delta z\sim0.3$ bin at intermediate redshifts, will test this prediction directly.
\end{abstract}

\section*{Introduction}

The physical nature of dark energy is unknown~\cite{riess1998,spergel2003}. The DESI collaboration finds evidence for dynamical dark energy at $2.8$--$4.2\sigma$~\cite{desi2025}. Existing models---quintessence~\cite{tsujikawa2013}, emergent dark energy~\cite{li2024}, and the CPL parameterization~\cite{chevallier2001}---describe the evolution of the equation of state $w(z)$ phenomenologically. None derives $w(z)$ from independently measured astrophysical quantities.

Here we test a hypothesis with that structure. Suppose cosmic structure formation occupies a finite reservoir of quantum distinguishability in spacetime---an idea motivated by holographic entropy bounds~\cite{thooft1993,susskind1995} and entropic gravity~\cite{verlinde2011,jacobson1995}. The dark energy density is then the energy of the remaining, unoccupied modes: $\rho_{\rm DE}(t)=\rho_\Lambda q(t)$, where $q(t)\in[0,1]$ is the fraction free. Crucially, $q(t)$ evolves because structure formation progressively occupies modes, and the equation of state $w(z)$ is determined by the star formation history---a quantity measured independently of dark energy.

This is a phenomenological model. The coupling between structure formation and mode occupation is calibrated from $w_0$, not derived from first principles. What the model provides is a specific, falsifiable shape for $w(z)$ that differs qualitatively from all existing parameterizations.

\section*{Model and self-consistent solution}

The occupied fraction $\Delta\equiv1-q$ obeys a nonlinear telegraph equation. On homogeneous cosmological scales:

\begin{equation}\label{eq:ode}
\gamma_0\Delta^n\frac{d\Delta}{dt} + \kappa\int_0^t\Delta\,dt' = \eta\,\dot{\rho}_*(t),
\end{equation}

where $\dot{\rho}_*(t)={\rm SFR}(t)$, $\gamma_0$ is a damping rate, $\kappa$ a memory strength, and $\eta$ a coupling. The dark energy density is $\rho_{\rm DE}=\rho_\Lambda(1-\Delta)$. In the regime where the driving term dominates ($z\gtrsim0.5$), the memory term is subdominant and equation (\ref{eq:ode}) yields $\Delta^{n+1}\propto\rho_*(t)$, with $\rho_*$ the cumulative stellar mass density. The data constrain $\kappa$ to be consistent with zero (68\% upper limit $\kappa<0.05$ in dimensionless units; Supplementary Information). From energy conservation,

\begin{equation}\label{eq:shape}
w(z)+1 = C\cdot\frac{{\rm SFR}(z)}{H(z)\,\rho_*(z)^{\,n/(n+1)}}\cdot\frac{1}{1-\Delta(z)},
\end{equation}

with $C\propto\sqrt{\eta/\gamma_0}$ calibrated from $w_0$. The natural value $n=1$ (damping proportional to deficit) is used; a scan over $n\in[0.3,3.0]$ changes $\chi^2$ by less than one unit.

Equation (\ref{eq:shape}) depends on $H(z)$ on both sides because $\rho_*(z)=\int_z^\infty{\rm SFR}(z')[(1+z')H(z')]^{-1}dz'$. We solve by Picard iteration: starting from $\Lambda$CDM $H(z)$, we compute $\rho_*$, $w$, and $\Delta$; update $H(z)$ via the Friedmann equation; recompute $\rho_*$; and repeat. Convergence (max$|\delta H/H|<10^{-5}$) occurs in $<15$ steps, driven by negative feedback: larger $H$ reduces $|dt/dz|$, reducing $\rho_*$, bringing $H$ back toward $\Lambda$CDM. We use the Madau \& Dickinson (2014) SFR fit~\cite{madau2014} and Planck 2018 parameters~\cite{planck2018}, with $r_d$ free. A robustness check using the Behroozi et al.\ (2019) SFR history shifts $(w_0,w_a)$ by $(+0.03,-0.08)$; the peak structure is preserved.

\section*{Comparison with DESI DR2}

We compare the converged prediction to DESI DR2 BAO distances~\cite{desi2025} at six effective redshifts (12 data points: $D_M/r_d$ and $D_H/r_d$ each), including the anti-correlation $\rho\approx-0.5$ between the two distance measures.

\begin{table}[htbp]
\centering
\caption{Model comparison against DESI DR2 BAO distances ($r_d$ free).}
\label{tab:models}
\begin{tabular}{lccccc}
\toprule
Model & DE params & $r_d$ (Mpc) & $\chi^2$ & AIC & $\Delta$AIC \\
\midrule
$\Lambda$CDM     & 0 ($w\equiv-1$)   & 148.6 & 15.0 & 17.0 & 0 \\
DGF ($\kappa=0$) & 2 ($C$, $n$)      & 146.0 & 19.1 & 25.1 & $+8.1$ \\
DGF ($\kappa$ free) & 3              & 146.0 & 19.1 & 27.1 & $+10.1$ \\
\bottomrule
\end{tabular}
\end{table}

Table~\ref{tab:models} reports the results. With $r_d$ free, $\Lambda$CDM achieves $\chi^2=15.0$ (best-fit $r_d=148.6$~Mpc). The DGF model yields $\chi^2=19.1$ and $r_d=146.0$~Mpc. The AIC difference of $+8.1$ favors $\Lambda$CDM. Adding the memory term as a free parameter further penalizes AIC without improving the fit.

The converged $w(z)$ is shown in Fig.~\ref{fig:wz}. The peak is mild: $w$ rises from $-0.80$ at $z=0$ to approximately $-0.31$ at $z\approx1.5$, then falls back toward $-1$ at higher redshift. The predicted Hubble parameter deviates from $\Lambda$CDM by approximately $2\%$ across all measured redshifts. The shaded band indicates the propagated $\pm27\%$ systematic uncertainty in the local star formation rate normalization.

\begin{figure}[htbp]
\centering
\includegraphics[width=0.85\textwidth]{../../LP32-S8_DESI-Contour/figures/fig2_wz_shape.png}
\caption{Converged $w(z)$ for $n=1$ (blue), with $n=0.7$ and $n=1.3$. The peak is mild---$w$ never crosses zero---due to negative feedback in the iteration. The band shows $\pm27\%$ SFR uncertainty.}
\label{fig:wz}
\end{figure}

\section*{Discussion}

The non-monotonic $w(z)$ is this model's distinguishing prediction. Table~\ref{tab:competitors} compares it to existing dark energy models. The CPL parameterization $w(a)=w_0+w_a(1-a)$ is linear in $(1-a)$ and cannot produce a peak. Quintessence with power-law potentials $V(\phi)\propto\phi^{-\alpha}$ yields monotonic $w(z)$~\cite{tsujikawa2013}. The emergent dark energy model~\cite{li2024} can produce non-monotonic behavior in some parameter regimes, but lacks a connection to the cosmic star formation history. The DGF model is unique in predicting a specific, SFR-determined peak.

\begin{table}[htbp]
\centering
\caption{Comparison of dark energy models.}
\label{tab:competitors}
\begin{tabular}{lccccc}
\toprule
Model & $N_{\rm par}$ & Physical origin? & Non-monotonic $w(z)$? & Unique prediction? \\
\midrule
$\Lambda$CDM & 0 & No & No & No \\
CPL & 2 & No & No & No \\
Quintessence ($V\propto\phi^{-\alpha}$) & 2--3 & Scalar field & No & No \\
Emergent DE & 2--3 & Integration constant & Possible & No \\
{\bf DGF} & {\bf 2} & {\bf Star formation} & {\bf Yes} & {\bf w(z) peak from SFR} \\
\bottomrule
\end{tabular}
\end{table}

The model's current limitations are significant and we state them explicitly. The coupling $\eta$ is calibrated from $w_0$, not derived. The star formation history input assumes $\Lambda$CDM luminosity distances; our iteration accounts for $H(z)$-dependence in $|dt/dz|$ but a fully self-consistent recomputation of the SFR from raw survey photometry is needed. The model is not statistically preferred over $\Lambda$CDM on current BAO data.

What the model provides is a specific, falsifiable prediction. The peak in $w(z)$ at intermediate redshift---a feature that the standard parameterizations cannot produce---will be tested by DESI DR3 and Euclid. If $w(z)$ is found to be monotonic, the model is excluded. If a peak is detected, it points toward a connection between dark energy and the cosmic history of structure formation.

\begin{thebibliography}{99}
\bibitem{riess1998} Riess, A.G.\ et al.\ {\it Astron.\ J.} {\bf 116}, 1009 (1998).
\bibitem{spergel2003} Spergel, D.N.\ et al.\ {\it Astrophys.\ J.\ Suppl.} {\bf 148}, 175 (2003).
\bibitem{desi2025} DESI Collaboration.\ {\it Phys.\ Rev.\ D} {\bf 112}, 083515 (2025).
\bibitem{chevallier2001} Chevallier, M.\ \& Polarski, D.\ {\it Int.\ J.\ Mod.\ Phys.\ D} {\bf 10}, 213 (2001).
\bibitem{li2026} Li, T.-N.\ et al.\ {\it Phys.\ Dark Universe} (2026); arXiv:2511.22512.
\bibitem{tsujikawa2013} Tsujikawa, S.\ {\it Class.\ Quant.\ Grav.} {\bf 30}, 214003 (2013).
\bibitem{li2024} Li, X.\ et al.\ {\it Phys.\ Rev.\ D} {\bf 109}, 043510 (2024).
\bibitem{thooft1993} 't Hooft, G.\ arXiv:gr-qc/9310026 (1993).
\bibitem{susskind1995} Susskind, L.\ {\it J.\ Math.\ Phys.} {\bf 36}, 6377 (1995).
\bibitem{verlinde2011} Verlinde, E.\ {\it J.\ High Energy Phys.} {\bf 2011}, 29 (2011).
\bibitem{jacobson1995} Jacobson, T.\ {\it Phys.\ Rev.\ Lett.} {\bf 75}, 1260 (1995).
\bibitem{madau2014} Madau, P.\ \& Dickinson, M.\ {\it Annu.\ Rev.\ Astron.\ Astrophys.} {\bf 52}, 415 (2014).
\bibitem{planck2018} Planck Collaboration.\ {\it Astron.\ Astrophys.} {\bf 641}, A6 (2020).
\bibitem{behroozi2019} Behroozi, P.\ et al.\ {\it Mon.\ Not.\ R.\ Astron.\ Soc.} {\bf 488}, 3143 (2019).
\end{thebibliography}

\end{document}
